% !TEX root = ../../main.tex
\setchapterimage{figure/P2/C1封面图.jpg}
\setchapterintro{以群、环、域为主线，掌握子结构、同态与商结构的关系。复习时先核对定义的条件，再理解每个构造为何成立，最后用例子与反例检验结论。}
\chapter{代数基本概念}
\label{p2:ch1}


\textbf{记号与约定：}
群的运算通常写成乘法，单位元记作$e$，元素的阶记作$\operatorname{ord}(a)$；加法群的单位元写成$0$。
子群用$H\leq G$表示（允许$H=G$），正规子群用$H\trianglelefteq G$表示。
环通常记为$R$。这里的环不必交换，也不必有乘法单位元；子环不要求与母环共用单位元；环同态不预先要求保持$1$。
这里的“理想”默认是双边理想。域和体均要求$1\neq0$。


\section{代数运算}


\begin{definition}
非空集合$A$上的一个\textbf{代数运算（二元运算）}是映射
\[
\mu:A\times A\longrightarrow A,\qquad(a,b)\longmapsto a*b.
\]
每一对有序元素都对应唯一的结果，且结果仍在$A$内。
\end{definition}

\begin{example}
整数加法、实数乘法、同阶方阵乘法、一个集合上的自映射复合都是代数运算。
实数上的$(x,y)\mapsto x$、$(x,y)\mapsto0$、$(x,y)\mapsto x+y^2$也是代数运算。
整数除法不是$\Z$上的代数运算：除数可以是$0$，且商未必是整数。
向量的数乘$\F\times V\to V$和内积$V\times V\to\F$一般不属于这里的二元运算。
\end{example}

\begin{definition}
对$A$上的运算$*$，若任意$a,b,c\in A$满足$(a*b)*c=a*(b*c)$，称它满足\textbf{结合律}。


若$a*b=b*a$，称它满足\textbf{交换律}。


若$e*a=a*e=a$，称$e$为\textbf{单位元}。

若在有单位元时$a*b=b*a=e$，称$b$是$a$的\textbf{逆元}。
\end{definition}

\textbf{注：}
封闭性和结果的唯一性已经包含在“映射$A\times A\to A$”中。
结合律允许改变括号，交换律允许改变元素的次序，两者不能混用。


\section{群的定义和简单性质}
\begin{definition}
非空集合$G$连同其上的乘法称为\textbf{群}，若满足：
\begin{enumerate}
\item $(ab)c=a(bc)$；
\item 存在$e\in G$，使$ea=a$对所有$a\in G$成立。称$e$为\textbf{左单位元}；
\item 每个$a\in G$存在$b\in G$，使$ba=e$，称$b$为\textbf{左逆元}。
\end{enumerate}


若另有$ab=ba$，则称为\textbf{交换群（阿贝尔群）}。
\end{definition}
下面我们可以证明，有左单位元，就有右单位元；有左逆元，就有右逆元。在此之前我们举几个例子。

\begin{proposition}
在上述公理下，$ba=e$必推出$ab=e$，且$ae=a$对所有$a$成立。
\end{proposition}
\begin{proof}
取$c$使$cb=e$。由结合律与左单位元性质，
\[
a=ea=(cb)a=c(ba)=ce,
\qquad ab=(ce)b=c(eb)=cb=e.
\]
再由$ba=e$和$ab=e$得$ae=a(ba)=(ab)a=ea=a$。
全过程没有预先使用右单位元或消去律。
\end{proof}

\begin{proposition}
群中单位元唯一，每个元素有唯一逆元，并且
\[
(a^{-1})^{-1}=a,\qquad e^{-1}=e,\qquad (ab)^{-1}=b^{-1}a^{-1}.
\]
\end{proposition}
\begin{proof}
若$e,f$均为单位元，则$e=ef=f$。
若$b,c$均为$a$的逆元，则$b=b(ac)=(ba)c=c$。
由于$a$是$a^{-1}$的逆元，得第一式；第二式直接验证。
最后，$(ab)(b^{-1}a^{-1})=a(bb^{-1})a^{-1}=e$，反向相乘也为$e$，由唯一性得第三式。
\end{proof}

\begin{proposition}
对任意$a,b\in G$，
\[
ax=b\iff x=a^{-1}b,\qquad ya=b\iff y=ba^{-1}.
\]
故两个方程均有唯一解，且$ab=ac\Rightarrow b=c$，$ba=ca\Rightarrow b=c$。
\end{proposition}
\begin{proof}
分别左乘、右乘$a^{-1}$得到解，代回即验证存在性。
对等式两侧左乘或右乘$a^{-1}$即得消去律。
\end{proof}



\textbf{记号：}
对$a\in G$与正整数$n$定义$a^n$为$n$个$a$的积，约定$a^0=e$，$a^{-n}=(a^{-1})^n$。
加法记号下对应整数倍$na$。


\begin{proposition}
\textbf{整数指数法则：}任意$m,n\in\Z$均有
\[
a^ma^n=a^{m+n},\qquad(a^m)^n=a^{mn}.
\]
若$ab=ba$，则还有$(ab)^n=a^nb^n$对所有$n\in\Z$成立。
\end{proposition}
\begin{proof}
非负指数的第一式由连乘积的结合律得到，第二式由分组得到。
负指数部分可用$a^{-1}$代替$a$；异号相乘时，把成对的$a,a^{-1}$消去，剩余部分正好给出指数之和。
同一元素的幂彼此交换，故负指数的第二式也由$(a^m)^{-1}=a^{-m}$得到。
若$ab=ba$，对$n\geq0$将$n$个$ab$中的$a$全部移至前面即可。
对$n<0$，用$(ab)^{-1}=b^{-1}a^{-1}=a^{-1}b^{-1}$归结到正指数。
\end{proof}

\begin{example}
\begin{enumerate}
	\item $(\Z,+),(\Q,+),(\R,+),(\C,+)$是交换群。
	\item $(\R\setminus\{0\},\cdot)$与$(\R_{>0},\cdot)$也是交换群。
	\item $U_n=\{z\in\C:z^n=1\}$在乘法下成$n$阶群：若$z,w\in U_n$，则$(zw)^n=1$，且$z^{-1}\in U_n$。
	\item 数域$P$上的$\GL_n(P)$、$\SL_n(P)$在矩阵乘法下成群。
	它们的封闭性和逆元条件分别由可逆性及$\det(AB)=\det A\det B$验证。
	\item 欧氏空间的正交变换也成群，因为复合与逆变换均保持内积。
\end{enumerate}
\end{example}
\begin{proposition}
	对交换群中的任意两个元素$a,b$及$n\in\Z$，有
	\[
	(ab)^n=a^n b^n.
	\]
\end{proposition}
\begin{proof}
	当$n\geq0$时，由交换律可将$n$个$ab$中的$a$移到前面，得到$(ab)^n=a^nb^n$。
	当$n<0$时，利用
	\[
	(ab)^{-1}=b^{-1}a^{-1}=a^{-1}b^{-1}
	\]
	归结为正整数情形。
\end{proof}
\begin{definition}
群的\textbf{阶}$|G|$是其元素个数。元素个数有限称\textbf{有限群}，否则称\textbf{无限群}。

\end{definition}

\section{群的例子}
\subsection*{图形的对称群与二面体群}
\begin{proposition}
平面图形$F$的所有保持$F$整体不变的平面等距变换，在复合下成群，称$F$的\textbf{对称群}。
\end{proposition}
\begin{proof}
恒等变换保持$F$。两个这样的变换的复合仍把$F$映为$F$，映射复合满足结合律。
若$f(F)=F$且$f$可逆，则$f^{-1}(F)=F$，故逆变换仍在该集合中。
\end{proof}

\begin{definition}
正$n$边形（$n\geq3$）的对称群称\textbf{二面体群}$D_n$，其阶为$2n$。
设$T$为绕中心旋转$2\pi/n$，$S$为对一条固定对称轴上的反射，则
\[
D_n=\{T^k,ST^k:0\leq k<n\},\qquad T^n=e,\quad S^2=e,\quad ST=T^{-1}S.
\]
\end{definition}

\begin{proposition}
上述$2n$个变换两两不同，且已穷尽正$n$边形的对称变换；当$n\geq3$时$D_n$不交换。
\end{proposition}
\begin{proof}
对称变换保持顶点集合，因而固定所有顶点的重心，即正多边形的中心。中心与两个相邻顶点不共线，而平面等距变换由三个不共线点的像唯一决定。
因此对称变换由一个顶点的像和相邻顶点的像唯一决定。

第一个顶点有$n$个选择，相邻顶点的像必须与它相邻，故有两个选择，至多有$2n$个对称变换。

$n$个旋转彼此不同，$n$个反射彼此不同；旋转保持平面定向而反射反转定向，故两类不重合。
关于过中心直线反射会反转旋转角，所以$STS=T^{-1}$。
若$ST=TS$，则$T=T^{-1}$，即$T^2=e$，与$T$的阶$n\geq3$矛盾。
\end{proof}

\begin{example}
正方形的对称群为$D_4=\{e,T,T^2,T^3,S,ST,ST^2,ST^3\}$，共有8个元素。
计算时先用$ST^k=T^{-k}S$把表达式化成$T^k$或$ST^k$，再将指数模4化简。
\end{example}

\subsection*{对称群与轮换}
\begin{definition}
集合$M$上的所有双射在复合下构成\textbf{全变换群}$S(M)$。
当$M=\{1,\ldots,n\}$时记作$S_n$，称\textbf{$n$元对称群}。
其元素称置换，记作
\[
\sigma=\begin{pmatrix}1&2&\cdots&n\\\sigma(1)&\sigma(2)&\cdots&\sigma(n)\end{pmatrix}.
\]
约定$(\sigma\tau)(i)=\sigma(\tau(i))$，即\textbf{右边先作用}。
\end{definition}

\begin{proposition}
$|S_n|=n!$。
\end{proposition}
\begin{proof}
双射的像序列恰是$1,\ldots,n$的一个排列，逐次选择各像有$n(n-1)\cdots1=n!$种。
\end{proof}

\begin{definition}
置换$(a_1\ a_2\ \cdots\ a_m)$把$a_1$映为$a_2$，依次轮换，最后把$a_m$映为$a_1$，并保持其余元素不变，称\textbf{轮换}。
长度为2的轮换称\textbf{对换}。两个轮换所涉及的元素集合不交时称\textbf{不相交}。
长度为1的轮换是恒等作用，通常省略。
\end{definition}

\begin{proposition}
不相交的轮换彼此交换。
\end{proposition}
\begin{proof}
对任意元素$x$分类：若$x$在其中一个轮换包含的元素中，另一个轮换同时固定$x$及其像；若$x$在两个包含的元素外，两者都固定$x$。
因此两种复合在每个$x$上的结果相同。对多个两两不相交的轮换，可依次应用这一结论。
\end{proof}

\begin{theorem}
\textbf{不相交轮换分解：}有限集合$\{1,\ldots,n\}$上的每个置换$\sigma\in S_n$都能写成不相交轮换的积。
分解在忽略轮换次序、每个轮换的循环起点以及省略的不动点后唯一。
\end{theorem}
\begin{proof}
对任意$x$连续作用$\sigma$。有限性保证出现重复，双射性保证第一次重复必回到$x$，否则与第一次重复的选取是矛盾的。

故$x,\sigma(x),\ldots,\sigma^{m-1}(x)$组成一个轮换。

从尚未出现的元素继续这一过程。若两次得到的轮换含有同一个元素，从该元素开始连续作用$\sigma$，就会依次得到相同的元素，直到回到起点，所以这两个轮换只可能是循环起点不同。因此新得到的轮换与此前的轮换不相交；重复这一过程，便能把全体元素分入这些轮换。

再证明唯一性。在任意不相交轮换分解中，取一个满足$\sigma(x)\neq x$的元素$x$，它必须出现在某个轮换中，且只能出现在一个轮换中。其余轮换都固定这个轮换中的元素，所以从$x$开始，这个轮换中后续的元素只能依次为$\sigma(x),\sigma^2(x),\ldots$，直到第一次回到$x$。因此，包含$x$的轮换完全由$\sigma$确定，只能改变循环起点。对其余被移动的元素重复这一论证，便确定了全部长度大于1的轮换；不动点对应长度为1的轮换，可以省略。因此分解唯一到上述意义。
\end{proof}

\begin{example}
\[
\begin{pmatrix}1&2&3&4&5&6\\3&1&2&5&4&6\end{pmatrix}=(132)(45).
\]
若$\sigma=(1243),\tau=(132)$，则$\sigma\tau=(34)$；须按右边先作用逐点检查。
\end{example}

\section{子群、陪集}


\begin{definition}
群$G$的非空子集$H$若在$G$的原运算下成群，称为\textbf{子群}，记为$H\leq G$。
全变换群的子群称\textbf{变换群}。
\end{definition}

\begin{definition}
\textbf{平凡子群与非平凡子群：}设$G$为群，单位元为$e$。子群$\{e\}$与$G$称为$G$的\textbf{平凡子群}。
若子群$H\leq G$满足$H\neq\{e\}$且$H\neq G$，即
\[
\{e\}\subsetneq H\subsetneq G,
\]
则称$H$为$G$的\textbf{非平凡子群}。
\end{definition}

\begin{theorem}
\textbf{子群判别法：}非空子集$H\subseteq G$为子群，当且仅当
\[
\forall a,b\in H,\qquad ab^{-1}\in H.
\]
等价地，$H$非空、对乘法封闭、对逆元封闭。
\end{theorem}
\begin{proof}
必要性由群公理成立。反之，取$a\in H$得$aa^{-1}=e\in H$。
由$e,b\in H$得$b^{-1}=eb^{-1}\in H$，再由$a,b^{-1}\in H$得$ab\in H$。
结合律从$G$继承，故$H$成群。
\end{proof}

\begin{proposition}
任意一族$G$的子群的交仍为子群（空族的交约定为$G$）。子群的并一般不是子群。
\end{proposition}
\begin{proof}
交集含$e$；其中$a,b$属于每个子群，故$ab^{-1}$也属于每个子群。

对于并集可以举反例：$2\Z\cup3\Z$含$2,3$却不含$5$，对加法不封闭。
\end{proof}

\subsection*{交错群}
\begin{definition}
令$D(x_1,\ldots,x_n)=\prod_{i<j}(x_i-x_j)$。
置换$\sigma$变量后，每一对差因子至多改变符号，故有唯一的$\operatorname{sgn}(\sigma)\in\{1,-1\}$使
\[
D(x_{\sigma(1)},\ldots,x_{\sigma(n)})=\operatorname{sgn}(\sigma)D(x_1,\ldots,x_n).
\]
符号为$1$称\textbf{偶置换}，为$-1$称\textbf{奇置换}。
全体偶置换记为$A_n$，称\textbf{交错群}。
\end{definition}

\begin{proposition}
$\operatorname{sgn}(\sigma\tau)=\operatorname{sgn}(\sigma)\operatorname{sgn}(\tau)$。
故$A_n\leq S_n$；当$n\geq2$时$|A_n|=n!/2$。
\end{proposition}
\begin{proof}
连续代换变量，把两个符号因子相乘即得第一式。

恒等置换符号为$1$，逆置换符号与原置换相同，因此子群判别法成立。

任意对换会使$D$变号（交换两变量时，二者对应的差变号，其余因子成对交换），故符号为$-1$。

固定一个对换$t$，映射$\sigma\mapsto t\sigma$是奇、偶置换集合间的双射，二者各占$n!$的一半。
当$n=1$时$A_1=S_1$，阶为1，不能用$n!/2$。

\end{proof}

\subsection*{陪集及其计数}
\begin{definition}
设$H\leq G$。$aH=\{ah:h\in H\}$称\textbf{左陪集}，$Ha=\{ha:h\in H\}$称\textbf{右陪集}。
$a$称陪集代表。左陪集个数记为$[G:H]$，称\textbf{指数}。
\end{definition}

\begin{proposition}
对$a,b\in G,H\leqslant G$，有
\[
aH=bH\iff a^{-1}b\in H,
\qquad Ha=Hb\iff ab^{-1}\in H.
\]
每个陪集中任一元素都可以作为\textbf{陪集代表}。任何两个左陪集（或两个右陪集）相等或不交。
\end{proposition}
\begin{proof}
若$aH=bH$，则$b=ah$，所以$a^{-1}b=h\in H$；反之$b=ah$给出$bH=ahH=aH$。
右陪集同理。

若$x\in aH\cap bH$，则$xH=aH$且$xH=bH$，故两个陪集相等。
\end{proof}

\begin{proposition}
$h\mapsto ah$给出$H\to aH$的双射，$h\mapsto ha$给出$H\to Ha$的双射。
左右陪集数相同，且$G$是所有不同左陪集的不交并，也是所有不同右陪集的不交并。
\end{proposition}
\begin{proof}
两个映射的逆分别是$x\mapsto a^{-1}x$和$x\mapsto xa^{-1}$。

每个$a$属于$aH$及$Ha$，所以陪集覆盖$G$。

映射$aH\mapsto Ha^{-1}$由取逆给出左右陪集集合的双射。
\end{proof}

\begin{theorem}
\textbf{拉格朗日定理：}若$G$有限且$H\leq G$，则
\[
|G|=[G:H]|H|,\qquad |H|\mid|G|.
\]
\end{theorem}
\begin{proof}
将$G$分为$[G:H]$个不交左陪集，每个陪集均有$|H|$个元素，相加即得。
\end{proof}

\begin{definition}
$\langle a\rangle=\{a^m:m\in\Z\}$称\textbf{由$a$生成的子群}。
$\langle a\rangle$的阶称\textbf{元素$a$的阶}，记为$\operatorname{ord}(a)$。
\end{definition}

\begin{proposition}
对$\langle a\rangle$，
若$a$的幂两两不同，$a$为无限阶元素；否则存在最小正整数$d$使$a^d=e$，且
\[
\langle a\rangle=\{e,a,\ldots,a^{d-1}\},\qquad \operatorname{ord}(a)=d,
\qquad a^m=e\iff d\mid m.
\]

$\langle a\rangle$是包含$a$的最小子群。
\end{proposition}
\begin{proof}
任何包含$a$的子群都包含其整数幂。
若$a^i=a^j$且$i<j$，则$a^{j-i}=e$。取最小正指数$d$。
对任意整数$m$作带余除法$m=qd+r$，$0\leq r<d$，有$a^m=a^r$。
最小性保证$e,a,\ldots,a^{d-1}$两两不同，也保证$a^r=e$当且仅当$r=0$。
\end{proof}

\begin{corollary}
若$|G|=n<\infty$，则每个元素的阶整除$n$，且$a^n=e$。
\end{corollary}
\begin{proof}
对$\langle a\rangle$应用拉格朗日定理，再写$n=\operatorname{ord}(a)q$。
\end{proof}

\textbf{注：}
拉格朗日定理只给出必要条件，不能据此断言每个$|G|$的因子都对应一个子群。
陪集一般只是子集，不是子群；含$e$的陪集才等于$H$。


\section{群的同构}
\begin{definition}
群$G,G'$之间的双射$\varphi$若满足$\varphi(ab)=\varphi(a)\varphi(b)$，称$\varphi$是\textbf{群同构}，且称为$G$同构于$G'$，记作$G\cong G'$。
\end{definition}

\begin{proposition}
同构保持单位元、逆元、元素的阶与交换性。
同构关系具有自反性、对称性、传递性。
\end{proposition}
\begin{proof}
$\varphi(e)\varphi(a)=\varphi(a)$，消去得$\varphi(e)=e'$；

$\varphi(a^{-1})\varphi(a)=e'$，故$\varphi(a^{-1})=\varphi(a)^{-1}$。

因此$\varphi(a^m)=\varphi(a)^m$，结合单射性得$a^m=e\iff\varphi(a)^m=e'$，阶保持。

对$ab=ba$应用$\varphi$即得像群交换；反向用$\varphi^{-1}$。

恒等映射是同构；若$u=\varphi(a),v=\varphi(b)$，则$\varphi^{-1}(uv)=ab$，所以逆映射是同构；两个同构的复合仍是保运算双射。
\end{proof}
\begin{theorem}
\textbf{凯莱定理：}每个群$G$都同构于集合$G$上的一个变换群。
\end{theorem}
\begin{proof}
设$G$是一个群。对每个$a\in G$，定义集合$G$到自身的变换
\[
L_a:G\longrightarrow G,\qquad L_a(x)=ax.
\]

先证明$L_a$可逆。对任意$x\in G$，有
\[
\begin{aligned}
(L_{a^{-1}}L_a)(x)
&=L_{a^{-1}}(ax)
=a^{-1}(ax)
=x,\\
(L_aL_{a^{-1}})(x)
&=L_a(a^{-1}x)
=a(a^{-1}x)
=x.
\end{aligned}
\]
因此两个复合都是恒等变换，即
\[
L_a^{-1}=L_{a^{-1}}.
\]
所以$L_a$是$G$上的双射。令
\[
G_l=\{L_a:a\in G\}\subseteq S(G).
\]

下面证明$G_l$是一个变换群。它非空，因为$L_e\in G_l$。
对任意$L_a,L_b\in G_l$及$x\in G$，有
\[(L_aL_b^{-1})(x)
=(L_aL_{b^{-1}})(x)=a(b^{-1}x)=(ab^{-1})x=L_{ab^{-1}}(x).\]
于是
\[
L_aL_b^{-1}=L_{ab^{-1}}\in G_l.
\]
由子群判别法，$G_l\leq S(G)$，所以$G_l$是集合$G$上的一个变换群。

最后证明$G$与$G_l$同构。定义映射
\[
\varphi:G\longrightarrow G_l,\qquad a\longmapsto L_a.
\]
由$G_l$的定义，$\varphi$是满射。又因为
\[
L_a(e)=ae=a,
\]
所以当$a\neq b$时，$L_a(e)\neq L_b(e)$，从而$L_a\neq L_b$。
因此$\varphi$也是单射，即$\varphi$是双射。

对任意$a,b,x\in G$，有
\[(L_aL_b)(x)=L_a(bx)=a(bx)=(ab)x=L_{ab}(x).\]
故$L_aL_b=L_{ab}$，即
\[
\varphi(ab)=L_{ab}=L_aL_b=\varphi(a)\varphi(b).
\]
所以$\varphi$是群同构，因而$G\cong G_l$，定理得证。
\end{proof}

\begin{proposition}
取右平移$R_a(x)=xa^{-1}$，则$R_aR_b=R_{ab}$；$a\mapsto R_a$也给出$G$到一个变换群的同构。
两种构造分别称\textbf{左正则表示、右正则表示}。
若$|G|=n$，它们均将$G$嵌入$S_n$。
\end{proposition}
\begin{proof}
$R_aR_b(x)=(xb^{-1})a^{-1}=x(ab)^{-1}=R_{ab}(x)$。
每个$R_a$可逆，且$R_a(e)=a^{-1}$保证单射。有限集合$G$的全变换群在选定编号后就是$S_n$。
\end{proof}

\textbf{注：}右正则表示中用$xa^{-1}$，而非$xa$，正是为了让$a\mapsto R_a$保持乘法次序。


\section{同态、正规子群}
\begin{definition}
映射$\varphi:G\to G'$满足$\varphi(ab)=\varphi(a)\varphi(b)$时称\textbf{群同态}。
像为$\varphi(G)=\{\varphi(a):a\in G\}$；核为
\[
\ker\varphi=\{a\in G:\varphi(a)=e'\}.
\]
单射同态称\textbf{嵌入}，满射同态称\textbf{满同态}。同构就是双射同态。
元素$y\in G'$的\textbf{完全反像}指集合$\varphi^{-1}(\{y\})$，不意味着同态有逆映射。
\end{definition}

\begin{proposition}
群同态满足$\varphi(e)=e'$，$\varphi(a^{-1})=\varphi(a)^{-1}$，$\varphi(a^m)=\varphi(a)^m$；
其像是$G'$的子群，核是$G$的子群；并且
\[
\varphi\text{为单射}\iff\ker\varphi=\{e\}.
\]
\end{proposition}
\begin{proof}
首先，由$ee=e$及同态性质，有
\[
\varphi(e)\varphi(e)=\varphi(ee)=\varphi(e).
\]
在$G'$中消去一个$\varphi(e)$，得$\varphi(e)=e'$。
再由
\[
\begin{aligned}
\varphi(a^{-1})\varphi(a)&=\varphi(a^{-1}a)=e',\\
\varphi(a)\varphi(a^{-1})&=\varphi(aa^{-1})=e',
\end{aligned}
\]
可知$\varphi(a^{-1})$是$\varphi(a)$的逆元，即
\[
\varphi(a^{-1})=\varphi(a)^{-1}.
\]
对正整数$m$反复使用保乘法性质，得
\[
\varphi(a^m)=\underbrace{\varphi(a)\cdots\varphi(a)}_{m\text{个}}=\varphi(a)^m.
\]
$m=0$时由单位元的结论成立；$m<0$时利用逆元的结论归结为正整数情形。

其次，像包含$\varphi(e)=e'$，故非空。若$\varphi(a),\varphi(b)$属于像，则
\[
\varphi(a)\varphi(b)^{-1}=\varphi(ab^{-1})\in\varphi(G).
\]
由子群判别法，$\varphi(G)\leq G'$。

对于核，$e\in\ker\varphi$。若$a,b\in\ker\varphi$，则
\[
\varphi(ab^{-1})=\varphi(a)\varphi(b)^{-1}=e'(e')^{-1}=e',
\]
所以$ab^{-1}\in\ker\varphi$。再次应用子群判别法，得到$\ker\varphi\leq G$。

最后，若$\ker\varphi=\{e\}$且$\varphi(a)=\varphi(b)$，则
\[
\varphi(a^{-1}b)=\varphi(a)^{-1}\varphi(b)=e',
\]
从而$a^{-1}b=e$，即$a=b$，所以$\varphi$为单射。
反之，若$\varphi$单射，对任意$a\in\ker\varphi$有
\[
\varphi(a)=e'=\varphi(e),
\]
故$a=e$，因此核恰为$\{e\}$。
\end{proof}

\begin{proposition}
\textbf{同态的纤维：}令$H=\ker\varphi$。若$y=\varphi(a)$，则
\[
\varphi^{-1}(\{y\})=aH=Ha.
\]
若$y\notin\varphi(G)$，完全反像为空。
\end{proposition}
\begin{proof}
若$x\in aH$，则$x=ah$对某个$h\in H$成立，于是
\[
\varphi(x)=\varphi(ah)=\varphi(a)\varphi(h)=\varphi(a)e'=y.
\]
因此$aH\subseteq\varphi^{-1}(\{y\})$。

反之，若$\varphi(x)=y=\varphi(a)$，则
\[
\varphi(a^{-1}x)=\varphi(a)^{-1}\varphi(x)=e',
\]
所以$a^{-1}x\in H$，即$x\in aH$。两侧包含合起来得到
\[
\varphi^{-1}(\{y\})=aH.
\]

右陪集可完全类似地证明。

若$y$不在像中，按像的定义不存在满足$\varphi(x)=y$的$x$，所以完全反像为空。
\end{proof}

\begin{definition}
若$H\leq G$且任意$a\in G$均有$aH=Ha$，称$H$为\textbf{正规子群}，记为$H\trianglelefteq G$。
\end{definition}

\begin{proposition}
对$H\leq G$，以下条件等价：
\[
H\trianglelefteq G;\qquad
\forall a\in G, aHa^{-1}=H;\qquad
\forall a\in G, aHa^{-1}\subseteq H.
\]
所有同态的核均正规；交换群的每个子群均正规。
\end{proposition}
\begin{proof}
对任意$a\in G$，在集合等式两侧右乘$a^{-1}$，有
\[
aH=Ha\iff aHa^{-1}=H.
\]
因此前两个条件等价，而共轭等式显然蕴含共轭包含关系。

反过来，设对每个$a\in G$都有$aHa^{-1}\subseteq H$。
将$a$替换为$a^{-1}$，得
\[
a^{-1}Ha\subseteq H.
\]
两侧再分别左乘$a$、右乘$a^{-1}$，得到
\[
H\subseteq aHa^{-1}.
\]
与原包含关系合并，便有$aHa^{-1}=H$，故三个条件等价。

若$H=\ker\varphi$，上一命题说明，对每个$a\in G$都有
\[
aH=\varphi^{-1}(\{\varphi(a)\})=Ha.
\]
所以核是正规子群。
若$G$交换，则$ah=ha$对任意$a\in G,h\in H$成立，因而同样有$aH=Ha$，所以它的每个子群都正规。
\end{proof}
\section{商群}
\begin{definition}
对$G$的子集$A,B$定义集合乘积$AB=\{ab:a\in A,b\in B\}$及$A^{-1}=\{a^{-1}:a\in A\}$。
\end{definition}
\begin{proposition}
集合乘积满足$(AB)C=A(BC)$，$(AB)^{-1}=B^{-1}A^{-1}$。
非空子集$H$为子群当且仅当$HH^{-1}\subseteq H$。
\end{proposition}
\begin{proof}[补充证明]
任取$x\in(AB)C$，存在$a\in A,b\in B,c\in C$使$x=(ab)c$。
由群的结合律，$x=a(bc)\in A(BC)$；反向包含同理，故
\[
(AB)C=A(BC).
\]

对任意$a\in A,b\in B$，由逆元法则有
\[
(ab)^{-1}=b^{-1}a^{-1}.
\]
遍历$a,b$便得到$(AB)^{-1}=B^{-1}A^{-1}$。

若$H$为子群，则$a,b\in H$时$ab^{-1}\in H$，所以$HH^{-1}\subseteq H$。
反过来，若$H$非空且$HH^{-1}\subseteq H$，则任意$a,b\in H$都有
\[
ab^{-1}\in HH^{-1}\subseteq H.
\]
由子群判别法，$H$为子群。
\end{proof}


\begin{theorem}
\textbf{陪集乘积与正规性：}子群$H\leq G$正规，当且仅当任意两个右陪集的集合乘积仍是一个右陪集；左陪集也有对应结论。
正规时$(Ha)(Hb)=Hab$。
\end{theorem}
\begin{proof}
\textbf{必要性：}设$H\trianglelefteq G$。对任意$a,b\in G$，由$aH=Ha$及$HH=H$，有
\[(Ha)(Hb)=H(aH)b=H(Ha)b=(HH)ab=Hab.\]
因此两个右陪集的集合乘积仍是右陪集。

\textbf{充分性：}假设任意两个右陪集的集合乘积仍是一个右陪集。
于是对任意$a,b\in G$，存在$c\in G$使
\[
(Ha)(Hb)=Hc.
\]
因为$e\in H$，所以$ab\in(Ha)(Hb)=Hc$，故$Hc=Hab$。因此
\[
HaHb=Hab.
\]
两侧右乘$b^{-1}$，得$HaH=Ha$。又因为$e\in H$，有
\[
aH\subseteq HaH=Ha,
\]
从而$aHa^{-1}\subseteq H$。

这一包含关系对每个$a\in G$都成立。把$a$换成$a^{-1}$，再在两侧左乘$a$、右乘$a^{-1}$，得到
\[
a^{-1}Ha\subseteq H
\quad\Longrightarrow\quad
H\subseteq aHa^{-1}.
\]
因此$aHa^{-1}=H$，故$H$正规。

左陪集的情形可由取逆得到。因为
\[
\bigl((aH)(bH)\bigr)^{-1}=(Hb^{-1})(Ha^{-1}),
\]
且左、右陪集在取逆下互相对应，所以左陪集乘积仍为左陪集，当且仅当右陪集有相应性质。
\end{proof}

\begin{definition}
若$H\trianglelefteq G$，全部不同陪集组成的集合记为$G/H$。
定义$(aH)(bH)=abH$，所得群称\textbf{商群}。
\end{definition}

\begin{proposition}
商群运算与代表选择无关；单位元为$H$，$aH$的逆元为$a^{-1}H$。
有限群有$|G/H|=|G|/|H|$。
\end{proposition}
\begin{proof}
由于$H$正规，左陪集与右陪集相同。前一定理给出
\[
(aH)(bH)=abH,
\]
其中左侧也恰是这两个集合的集合乘积。因此，若$a'H=aH$且$b'H=bH$，则
\[
a'b'H=(a'H)(b'H)=(aH)(bH)=abH.
\]
这证明运算不依赖所选的代表元素。

对任意$aH,bH,cH$，由$G$的结合律，有
\[
\begin{aligned}
\bigl((aH)(bH)\bigr)(cH)&=(ab)cH,\\
(aH)\bigl((bH)(cH)\bigr)&=a(bc)H.
\end{aligned}
\]
因为$(ab)c=a(bc)$，所以商群乘法满足结合律。

又有
\[
H(aH)=(aH)H=aH,
\]
因此$H=eH$为单位元。并且
\[
(aH)(a^{-1}H)=(a^{-1}H)(aH)=H,
\]
所以$aH$的逆元为$a^{-1}H$，这些陪集确实构成群。

当$G$有限时，$G/H$的元素个数就是不同陪集的个数。由拉格朗日定理，
\[
|G|=[G:H]|H|=|G/H||H|,
\]
所以$|G/H|=|G|/|H|$。
\end{proof}

\begin{definition}
\textbf{自然同态：}$\pi:G\to G/H$，$\pi(a)=aH$是满同态，$\ker\pi=H$。
因此$G$的正规子群恰是以$G$为定义域的群同态的核。
\end{definition}
\begin{theorem}\label{p2:group-first-iso}
\textbf{群同态基本定理（第一同构定理）：}任意同态$\varphi:G\to G'$均诱导同构
\[
\bar\varphi:G/\ker\varphi\longrightarrow\varphi(G),\qquad
a\ker\varphi\longmapsto\varphi(a).
\]
特别地，满同态给出$G/\ker\varphi\cong G'$。
\end{theorem}
\begin{proof}
令$N=\ker\varphi$，由上一节可知$N\trianglelefteq G$。
定义
\[
\bar\varphi:G/N\longrightarrow\varphi(G),\qquad
\bar\varphi(aN)=\varphi(a).
\]

先说明这个定义与代表的选择无关。上一节关于完全反像的结论给出
\[aN=bN\iff a^{-1}b\in N\iff\varphi(a^{-1}b)=e'\iff\varphi(a)^{-1}\varphi(b)=e'\iff\varphi(a)=\varphi(b).\]
从前往后看，说明同一个陪集不会对应不同的像，故映射良定义；从后往前看，说明像相同的两个陪集必相同，故映射单射。

任意$y\in\varphi(G)$都可写为$y=\varphi(a)$，所以
\[
y=\bar\varphi(aN),
\]
故映射满到$\varphi(G)$，从而是双射。

最后，对任意$aN,bN\in G/N$，有
\[\bar\varphi\bigl((aN)(bN)\bigr)=\bar\varphi(abN)=\varphi(ab)=\varphi(a)\varphi(b)=\bar\varphi(aN)\bar\varphi(bN).\]
因此$\bar\varphi$为群同构。若$\varphi$满，则$\varphi(G)=G'$，所以$G/N\cong G'$。
\end{proof}

\begin{definition}
以对象为顶点、映射为箭头组成的图，若起点与终点相同的各条路径对应相同的复合映射，称\textbf{交换图}。
上述定理给出
\[
\begin{tikzcd}
G \arrow[r,"\varphi"] \arrow[d,"\pi"'] & \varphi(G)\\
G/N \arrow[ur,"\bar\varphi"'] &
\end{tikzcd}
\qquad\varphi=\bar\varphi\circ\pi.
\]
\end{definition}

\begin{example}
行列式满同态给出$\GL_n(P)/\SL_n(P)\cong P^\times$。

符号满同态给出$S_n/A_n\cong\{1,-1\}$（$n\geq2$）。
\end{example}

\section{环、子环}
\begin{definition}
非空集合$R$带有加法和乘法，若$(R,+)$为交换群，乘法满足结合律，且两侧分配律成立：
\[
a(b+c)=ab+ac,\qquad(b+c)a=ba+ca,
\]
则称$R$为\textbf{环}。加法单位元记为$0$，加法逆元记为$-a$。
\end{definition}



\begin{proposition}
对$a,b\in R$，
\[
a0=0a=0,\quad(-a)b=a(-b)=-(ab),\quad(-a)(-b)=ab.
\]
\end{proposition}
\begin{proof}
取任意$b\in R$，有
\[
ab=a(b+0)=ab+a0.
\]
在加法群中消去$ab$，得$a0=0$。另一侧同样有
\[
ba=(b+0)a=ba+0a,
\]
消去$ba$，得到$0a=0$。

再证明符号法则。由分配律及刚得到的零元性质，
\[
\begin{aligned}
(-a)b+ab&=(-a+a)b=0b=0,\\
a(-b)+ab&=a(-b+b)=a0=0.
\end{aligned}
\]
因此$(-a)b$与$a(-b)$都是$ab$的加法逆元。由加法逆元的唯一性，
\[
(-a)b=a(-b)=-(ab).
\]
最后再次应用这一结论，得到
\[
\begin{aligned}
(-a)(-b)
&=-\bigl(a(-b)\bigr)\\
&=-\bigl(-(ab)\bigr)\\
&=ab.
\end{aligned}
\]
\end{proof}

\begin{proposition}
\textbf{整数倍与正整数幂：}定义$na$为加法群中的整数倍。对$m,n\in\Z$有
\[
(m+n)a=ma+na,\quad m(na)=(mn)a,\quad n(a+b)=na+nb,
\]
\[
(na)b=n(ab)=a(nb).
\]
定义正整数幂后，$a^ma^n=a^{m+n}$，$(a^m)^n=a^{mn}$对$m,n\geq1$成立。
\end{proposition}
\begin{proof}
因为$(R,+)$为交换群，把前面群的整数指数法则改用加法记号，就得到
\[
(m+n)a=ma+na,\qquad m(na)=(mn)a.
\]
又因为加法交换，同一整数倍可以分组，故
\[
n(a+b)=na+nb\qquad(n\in\Z).
\]

对正整数$n$，连续应用分配律，有
\[
\begin{aligned}
(na)b
&=(\underbrace{a+\cdots+a}_{n\text{个}})b\\
&=\underbrace{ab+\cdots+ab}_{n\text{个}}\\
&=n(ab),
\end{aligned}
\]
同样有$a(nb)=n(ab)$。当$n=0$时，这些等式由零元性质成立。
若$n=-k<0$，则由符号法则，
\[
\begin{aligned}
(na)b&=(-(ka))b=-\bigl((ka)b\bigr)=-k(ab)=n(ab),\\
a(nb)&=a(-(kb))=-\bigl(a(kb)\bigr)=-k(ab)=n(ab).
\end{aligned}
\]
因此$(na)b=n(ab)=a(nb)$对全部整数$n$成立。

最后，乘法结合律保证有限连乘积可以任意改变括号。对$m,n\geq1$，连接及分组相同的因子，得到
\[
\begin{aligned}
a^ma^n
&=(\underbrace{a\cdots a}_{m\text{个}})(\underbrace{a\cdots a}_{n\text{个}})
=a^{m+n},\\
(a^m)^n
&=\underbrace{(a^m)\cdots(a^m)}_{n\text{个}}
=a^{mn}.
\end{aligned}
\]
这里仅使用正整数幂，没有假定环有乘法单位元。
\end{proof}
\textbf{注：}
一般环中不能定义$a^0$，因为可能没有$1$；即使有$1$，非可逆元素也不能定义负整数幂。
整数倍$na$的左边$n$是整数，不能在无单位元的环中擅自把它看成环元素$n\cdot1$。


\begin{definition}
非空子集$S\subseteq R$若在$R$的两种运算下成环，称\textbf{子环}。
\end{definition}
\begin{theorem}
\textbf{子环判别法：}非空$S\subseteq R$为子环，当且仅当任意$a,b\in S$均有$a-b\in S$和$ab\in S$。
\end{theorem}
\begin{proof}
\textbf{必要性：}若$S$是子环，则$(S,+)$是$(R,+)$的子群，所以
\[
a,b\in S\Longrightarrow a-b\in S.
\]
又因为乘法是$S$上的运算，所以$ab\in S$。

\textbf{充分性：}设$S$非空，并且对差与积封闭。取$c\in S$，有
\[
0=c-c\in S.
\]
对$a,b\in S$，依次得到
\[
-a=0-a\in S,\qquad a+b=a-(-b)\in S.
\]
因此$S$在加法下是子群，且加法交换律从$R$继承。
乘积封闭保证乘法限制到$S$后仍是二元运算；乘法结合律及两侧分配律也从$R$继承，所以$S$成环。
\end{proof}

\begin{example}
$\Z\subseteq\Q$、$n\Z\subseteq\Z$均为子环。

$M_n(P)$中的对角矩阵、数量矩阵、上三角矩阵和下三角矩阵分别构成子环。

差和乘积保持对应的矩阵形状，即可用判别法验证。

映射$A\mapsto\operatorname{diag}(A,0)$把$M_2(P)$同构到$M_3(P)$的一个子环。
该子环的单位元是$\operatorname{diag}(1,1,0)$，与$M_3(P)$的单位元不同。
\end{example}

\begin{definition}
双射$\varphi:R\to R'$若同时保持加法和乘法，称\textbf{环同构}，记作$R\cong R'$。
\end{definition}
\begin{proposition}
环同构保持$0$、负元、整数倍、正整数幂；若母环有单位元，则同构将其映为目标环的单位元。
环同构关系具有自反性、对称性和传递性。
\end{proposition}
\begin{proof}[补充证明]
设$\varphi:R\to R'$为环同构。它是加法群的同构，所以
\[
\varphi(0)=0,
\qquad \varphi(-a)=-\varphi(a),
\qquad \varphi(na)=n\varphi(a)\quad(n\in\Z).
\]
由保持乘法，对正整数$m$反复应用同态性质，得到
\[
\varphi(a^m)=\varphi(a)^m.
\]

若$R$有单位元$1_R$，任取$y\in R'$，可写成$y=\varphi(a)$。
因此
\[
\varphi(1_R)y=\varphi(1_Ra)=y,
\qquad y\varphi(1_R)=\varphi(a1_R)=y.
\]
故$\varphi(1_R)$是$R'$的单位元。

恒等映射是环同构；环同构的逆映射仍保持加法与乘法；两个环同构的复合也保持这两种运算。
它们都是双射，所以同构关系具有自反性、对称性和传递性。
\end{proof}
\section{各种特殊类型的环}
\begin{definition}
有乘法单位元的环称\textbf{幺环}，单位元记为$1$。

幺环中有双边乘法逆元的元素称\textbf{单位（可逆元）}，全体单位记为$R^\times$。

乘法交换的环称\textbf{交换环}。
\end{definition}

\begin{proposition}
乘法单位元唯一；双边逆元唯一；若$a$有左逆$b$和右逆$c$，则$b=c$。
全体单位在乘法下成群，称\textbf{单位群}。
\end{proposition}
\begin{proof}
若$u,v$均为乘法单位元，则
\[
u=uv=v,
\]
故单位元唯一。若$ba=1$且$ac=1$，则
\[b=b(ac)=(ba)c=c.\]
所以左逆与右逆同时存在时必相等，双边逆元也因此唯一。

下面验证$R^\times$成群。它包含$1$，因为$1\cdot1=1$。
若$u,v\in R^\times$，则
\[
\begin{aligned}
(uv)(v^{-1}u^{-1})&=u(vv^{-1})u^{-1}=1,\\
(v^{-1}u^{-1})(uv)&=v^{-1}(u^{-1}u)v=1.
\end{aligned}
\]
故$uv$仍可逆，而且$(uv)^{-1}=v^{-1}u^{-1}$，说明单位集合对乘法封闭。
又因为$u^{-1}$的逆是$u$，所以每个单位的逆仍是单位。
结合律从环乘法继承，因此$R^\times$在乘法下成群。
\end{proof}

\begin{definition}
非零元素$a$若存在非零$b$使$ab=0$，称$a$为\textbf{左零因子}；
若存在非零$b$使$ba=0$，称$a$为\textbf{右零因子}。两者统称零因子。
无零因子的交换幺环，且$1\neq0$，称\textbf{整环}。
\end{definition}

\begin{proposition}
环无零因子时，非零元素可左右消去：$ab=ac,a\neq0\Rightarrow b=c$，$ba=ca,a\neq0\Rightarrow b=c$。
反之，若对所有非零元素两侧消去律均成立，则环无零因子。
\end{proposition}
\begin{proof}
若$a\neq0$且$ab=ac$，则
\[
\begin{aligned}
ab-ac&=0,\\
a(b-c)&=0.
\end{aligned}
\]
如果$b-c\neq0$，这就给出两个非零元素的乘积为0，与无零因子矛盾。因此$b-c=0$，即$b=c$。

右侧消去同理：若$ba=ca$，则
\[
(b-c)a=0.
\]
由$a\neq0$及无零因子，同样得到$b=c$。

\textit{补充逆命题。}假设非零元素的消去律成立。若$a,b\neq0$且$ab=0$，则
\[
ab=a0.
\]
因为$a\neq0$，左侧消去得$b=0$，矛盾。因此环中不存在非零元素的零乘积，也就没有零因子。
\end{proof}

\begin{example}

$\Z$是整环，但不是域。

$m\Z$（$m\geq2$）没有乘法单位元，所以不是整环。

若单位元$u\in m\Z$存在，由$u\cdot m=m$在整数中得$u=1$，矛盾。

$M_n(P)$（$n\geq2$）有零因子：$E_{11}E_{22}=0$，两者均非零。

幺环中的$1=0$等价于该环只有一个元素，因为$a=a1=a0=0$。
\end{example}

\begin{definition}

交换幺环$F$若$1\neq0$且每个非零元素可逆，称\textbf{域}。

去掉乘法交换条件，得到\textbf{体（除环）}。

域的子环若本身为域，称\textbf{子域}。
\end{definition}

\begin{proposition}
体无零因子，域是整环；但整环未必是域。
域的子域与母域具有相同的单位元。
\end{proposition}
\begin{proof}
若在体中$ab=0$且$a,b\neq0$，则$a^{-1}$存在，故
\[b=(a^{-1}a)b=a^{-1}(ab)=a^{-1}0=0.\]
矛盾，所以两个非零元素的乘积不可能为0，体无零因子。
域还是交换幺环且$1\neq0$，因此域是整环。
但$\Z$是整环，元素2在其中没有乘法逆元，所以整环未必为域。

设$K$是域$F$的子域，其单位元为$u$。
由于$K$的加法是$F$加法的限制，加法单位元相同；又因$K$为域，所以$u\neq0$。
由$u$在$K$中为单位元，得到$u^2=u$。在$F$中，$u$有逆元，故
\[
u^2=u
\quad\Longrightarrow\quad
u^{-1}u^2=u^{-1}u
\quad\Longrightarrow\quad
u=1.
\]
因此子域与母域的乘法单位元相同。
\end{proof}

\begin{theorem}
\textbf{有限整环是域。}
\end{theorem}
\begin{proof}
设
\[
R=\{a_1,a_2,\ldots,a_n\},\qquad a_1=1.
\]
任取非零$a\in R$，考虑$n$个乘积
\[
aa_1,aa_2,\ldots,aa_n.
\]
如果$aa_i=aa_j$，由整环中非零元素的消去律可知$a_i=a_j$，于是$i=j$。
因此这些乘积两两不同。

它们都属于只有$n$个元素的集合$R$，所以恰好列出了$R$的全部元素，其中必有单位元1。
故存在$k$使
\[
aa_k=1.
\]
整环的乘法交换，因此还有$a_ka=1$，所以$a_k$是$a$的双边逆元。
因为$a$是任取的非零元素，所有非零元素都可逆，所以$R$为域。
\end{proof}

\subsection*{四元数体的矩阵模型}
\begin{proposition}
集合
\[
\mathbb H=\left\{X(\alpha,\beta)=
\begin{pmatrix}\alpha&\beta\\-\bar\beta&\bar\alpha\end{pmatrix}:
\alpha,\beta\in\C\right\}
\]
是$M_2(\C)$的子环，并且是非交换的体，称\textbf{四元数体}。
\end{proposition}
\begin{proof}

零矩阵为$X(0,0)$，而且
\[
X(\alpha,\beta)-X(\gamma,\delta)
=X(\alpha-\gamma,\beta-\delta),
\]
所以对差封闭。乘法方面，直接计算得
\[
X(\alpha,\beta)X(\gamma,\delta)
=X(\alpha\gamma-\beta\bar\delta,\alpha\delta+\beta\bar\gamma),
\]
故乘法也封闭，由子环判别法知$\mathbb H$是$M_2(\C)$的子环。

取单位矩阵$X(1,0)$作为乘法单位元。
对$X=X(\alpha,\beta)\neq0$，$\alpha,\beta$不全为0，所以
\[
\det X=\alpha\bar\alpha+\beta\bar\beta=|\alpha|^2+|\beta|^2>0.
\]
因此$X$在矩阵环中可逆，逆矩阵为
\[
X^{-1}=\frac1{|\alpha|^2+|\beta|^2}
\begin{pmatrix}\bar\alpha&-\beta\\\bar\beta&\alpha\end{pmatrix}\in\mathbb H
\]
更具体地，记$d_0=|\alpha|^2+|\beta|^2$，则
\[
X^{-1}=X\left(\frac{\bar\alpha}{d_0},-\frac{\beta}{d_0}\right)\in\mathbb H.
\]
所以$\mathbb H$的每个非零元素都具有仍在$\mathbb H$中的乘法逆元，故$\mathbb H$是体。

最后把元素写成四个实系数的形式。令
\[
I=\begin{pmatrix}i&0\\0&-i\end{pmatrix},\quad
J=\begin{pmatrix}0&1\\-1&0\end{pmatrix},\quad
K=\begin{pmatrix}0&i\\i&0\end{pmatrix}.
\]
若$\alpha=a+bi,\beta=c+di$，则
\[
X(\alpha,\beta)
=\begin{pmatrix}a+bi&c+di\\-c+di&a-bi\end{pmatrix}
=a+bI+cJ+dK.
\]
复数的实部与虚部唯一，所以这四个实系数也唯一。
由矩阵乘法计算，得到
\[
I^2=J^2=K^2=-1,\quad IJ=K=-JI,\quad JK=I=-KJ,\quad KI=J=-IK.
\]
$IJ\neq JI$，所以不交换，不是域。
\end{proof}

\section{环的同态、理想}
\begin{definition}
映射$\varphi:R\to R'$若满足
\[
\varphi(a+b)=\varphi(a)+\varphi(b),\qquad\varphi(ab)=\varphi(a)\varphi(b),
\]
称\textbf{环同态}。像为$\varphi(R)$，核为$\ker\varphi=\{a:\varphi(a)=0\}$。
像为$\{0\}$称\textbf{零同态}；像为$R'$称\textbf{满同态}。
\end{definition}

\begin{proposition}
环同态满足$\varphi(0)=0$、$\varphi(-a)=-\varphi(a)$、$\varphi(na)=n\varphi(a)$。
像是子环，核是加法子群，且$\varphi$为单射当且仅当核为$\{0\}$。
\end{proposition}
\begin{proof}
先把环同态看成加法群同态。由前面的群同态性质，立即得到
\[
\varphi(0)=0,\qquad
\varphi(-a)=-\varphi(a),\qquad
\varphi(na)=n\varphi(a).
\]

像包含$\varphi(0)=0$，因此非空。
对像中的任意两个元素$\varphi(a),\varphi(b)$，有
\[
\begin{aligned}
\varphi(a)-\varphi(b)&=\varphi(a-b)\in\varphi(R),\\
\varphi(a)\varphi(b)&=\varphi(ab)\in\varphi(R).
\end{aligned}
\]
由子环判别法，$\varphi(R)$为$R'$的子环。

核就是加法群同态的核，所以它是$(R,+)$的子群。
又有
\[
\begin{aligned}
\varphi(a)=\varphi(b)
&\iff\varphi(a)-\varphi(b)=0\\
&\iff\varphi(a-b)=0\\
&\iff a-b\in\ker\varphi.
\end{aligned}
\]
若核为$\{0\}$，上式推出$a=b$，所以同态单射；反之，若同态单射，则
$\varphi(a)=0=\varphi(0)$只能在$a=0$时成立，故核为$\{0\}$。
\end{proof}

\textbf{注：}
依本章约定，环同态未必把母环的$1$映为目标环的$1$。
若定义域有$1$，$\varphi(1)$只是\textbf{像子环}的单位元；若同态还满且目标为幺环，则必有$\varphi(1)=1$。

\begin{definition}
环$R$的加法子群$I$若对所有$r\in R,a\in I$有$ra\in I$且$ar\in I$，称\textbf{双边理想}。
只要求$ra\in I$称\textbf{左理想}，只要求$ar\in I$称\textbf{右理想}。
$\{0\}$与$R$称平凡理想。
\end{definition}

\begin{proposition}
同态的核是双边理想。每个双边理想都是子环，但子环未必是理想。
\end{proposition}
\begin{proof}
先由加法群同态的性质知道$\ker\varphi$是加法子群。
取$a\in\ker\varphi$和任意$r\in R$，有
\[
\begin{aligned}
\varphi(ra)&=\varphi(r)\varphi(a)=\varphi(r)0=0,\\
\varphi(ar)&=\varphi(a)\varphi(r)=0\varphi(r)=0.
\end{aligned}
\]
因此$ra,ar\in\ker\varphi$，满足\textbf{双边吸收条件}，所以核是双边理想。

\textbf{补充理想与子环的比较：}若$I$为双边理想，则它作为加法子群非空，且
\[
a,b\in I\Longrightarrow a-b\in I.
\]
又因为$a\in I\subseteq R,b\in I$，由吸收性有$ab\in I$。
所以$I$对差与积封闭，由子环判别法为子环。

反过来，$\Z$是$\Q$的子环，但
\[
1\in\Z,\qquad \frac12\in\Q,\qquad \frac12\cdot1\notin\Z.
\]
因此$\Z$不满足被母环元素乘后仍在其中的条件，不是$\Q$的理想。
\end{proof}

\begin{proposition}
幺环$R$的理想$I$若含$1$，则$I=R$；若含任一单位，也有$I=R$。
\end{proposition}
\begin{proof}
若$1\in I$，任取$r\in R$，由理想的吸收性，
\[
r=r1\in I.
\]
因此$R\subseteq I$，而$I\subseteq R$，故$I=R$。

若$I$含任一可逆元$u$，则$u^{-1}\in R$，再次由吸收性，
\[
u^{-1}u=1\in I.
\]
于是归结为上一种情形，仍有$I=R$。
\end{proof}
\section{商环}
\begin{definition}
若$I$是$R$的双边理想，令$R/I=\{a+I:a\in R\}$。
定义
\[
(a+I)+(b+I)=a+b+I,\qquad(a+I)(b+I)=ab+I,
\]
所得环称\textbf{商环}。
\end{definition}

\begin{proposition}
商环运算良定义，且满足环公理。若$R$是幺环，其商环的单位元为$1+I$；若$I=R$，商环为零环。
\end{proposition}
\begin{proof}
先把$I$看成加法群的子群，使用加法商群定义陪集的加法。
因为加法交换，$I$在加法群中正规，所以全体$a+I$在加法下为交换群，零元为$I$，$a+I$的负元为$-a+I$。

再考虑乘法。取任意
\[
x=a+i\in a+I,\qquad y=b+j\in b+I,
\qquad i,j\in I.
\]
由分配律，
\[
\begin{aligned}
xy
&=(a+i)(b+j)\\
&=ab+ib+aj+ij.
\end{aligned}
\]
因为$I$是双边理想，$ib,aj,ij\in I$，所以$xy\in ab+I$。
这说明两个陪集的任意代表相乘，都落在同一个陪集$ab+I$中。
因此规定
\[
(a+I)(b+I)=ab+I
\]
与代表的选择无关，乘法良定义。

由$R$中的对应性质，
\[
\begin{aligned}
\bigl((a+I)(b+I)\bigr)(c+I)&=(ab)c+I,\\
(a+I)\bigl((b+I)(c+I)\bigr)&=a(bc)+I.
\end{aligned}
\]
两式相等，故乘法结合。又有
\[
\begin{aligned}
(a+I)\bigl((b+I)+(c+I)\bigr)
&=a(b+c)+I\\
&=ab+ac+I\\
&=(a+I)(b+I)+(a+I)(c+I),
\end{aligned}
\]
以及
\[
\begin{aligned}
\bigl((b+I)+(c+I)\bigr)(a+I)
&=(b+c)a+I\\
&=ba+ca+I\\
&=(b+I)(a+I)+(c+I)(a+I).
\end{aligned}
\]
所以两侧分配律成立，所得结构确为环。

若$R$有单位元1，则
\[
(1+I)(a+I)=a+I=(a+I)(1+I),
\]
故$1+I$为商环的单位元。若$I=R$，所有陪集都等于$R$，所以商环只有一个元素，即零环。
\end{proof}

\textbf{注：}
商环乘法是\textbf{定义}为$ab+I$。两个陪集作为集合的乘积一般仅有
$\{xy:x\in a+I,y\in b+I\}\subseteq ab+I$，未必相等。
例如$R=\Z,I=2\Z,a=b=0$，集合乘积为$4\Z$，而商环乘积对应陪集$2\Z$。

\begin{proposition}
自然映射$\pi:R\to R/I$，$a\mapsto a+I$是满环同态，且核为$I$。
所以双边理想恰是环同态的核。
\end{proposition}
\begin{proof}
以下直接按定义验证。
对任意$a,b\in R$，有
\[
\begin{aligned}
\pi(a+b)&=a+b+I=(a+I)+(b+I)=\pi(a)+\pi(b),\\
\pi(ab)&=ab+I=(a+I)(b+I)=\pi(a)\pi(b).
\end{aligned}
\]
所以$\pi$为环同态。
每个商环元素都是某个$a+I=\pi(a)$，故它是满射。

商环的零元为$I$，于是
\[
\begin{aligned}
\ker\pi
&=\{a\in R:a+I=I\}\\
&=I.
\end{aligned}
\]
因此每个双边理想都是一个环同态的核；反过来，同态的核已证明为双边理想。
\end{proof}

\begin{proposition}
对$n\geq2$，$\bar a\in\Z/n\Z$可逆当且仅当$\gcd(a,n)=1$。
其单位群(记为$(\Z/n\Z)^{\times}$)的阶为欧拉函数$\varphi(n)$。
$\Z/n\Z$是域当且仅当$n$为素数。
\end{proposition}
\begin{proof}
若$\bar a$可逆，则存在$\bar b$使
\[
\bar a\bar b=\bar1,
\]
即$ab\equiv1\pmod n$。因此存在$k\in\Z$使
\[
ab+kn=1.
\]
$a,n$的任一公因子都整除右边的1，故$\gcd(a,n)=1$。

反之，若$\gcd(a,n)=1$，由贝祖等式存在$b,k\in\Z$使$ab+kn=1$。
模$n$化简，得到
\[
\bar a\bar b=\overline{ab}=\bar1.
\]
因为商环乘法交换，$\bar b$是$\bar a$的双边逆元。
所以可逆元恰为模$n$的互素剩余类，个数为$\varphi(n)$。

当$n=p$为素数时，每个非零剩余类可取代表$1\leq a<p$，从而$\gcd(a,p)=1$，所以均可逆。
又因$p\geq2$，$\bar1\neq\bar0$，故$\Z/p\Z$为域。

若$n$为合数，可写为$n=uv$，其中$1<u,v<n$。
于是
\[
\bar u\neq\bar0,\qquad \bar v\neq\bar0,
\qquad \bar u\bar v=\overline{uv}=\bar n=\bar0.
\]
这给出零因子，所以$\Z/n\Z$不可能为域。
\end{proof}

\begin{theorem}\label{p2:ring-first-iso}
\textbf{环同态基本定理（第一同构定理）：}任意环同态$\varphi:R\to R'$诱导环同构
\[
R/\ker\varphi\cong\varphi(R),\qquad a+\ker\varphi\longmapsto\varphi(a).
\]
特别地，若$\varphi$满，则$R/\ker\varphi\cong R'$。
\end{theorem}
\begin{proof}
令$I=\ker\varphi$，它已证明为双边理想，故可作商环$R/I$。
定义
\[
\bar\varphi:R/I\longrightarrow\varphi(R),\qquad
\bar\varphi(a+I)=\varphi(a).
\]

先核对代表选择。由加法群的陪集性质，
\[
\begin{aligned}
a+I=b+I
&\iff a-b\in I\\
&\iff\varphi(a-b)=0\\
&\iff\varphi(a)-\varphi(b)=0\\
&\iff\varphi(a)=\varphi(b).
\end{aligned}
\]
正向蕴含说明$\bar\varphi$良定义，反向蕴含说明它为单射。
任意$y\in\varphi(R)$都有$y=\varphi(a)=\bar\varphi(a+I)$，所以它满到$\varphi(R)$，从而是双射。

再验证两种运算。对任意$a+I,b+I$，有
\[
\begin{aligned}
\bar\varphi\bigl((a+I)+(b+I)\bigr)
&=\bar\varphi(a+b+I)\\
&=\varphi(a+b)\\
&=\varphi(a)+\varphi(b)\\
&=\bar\varphi(a+I)+\bar\varphi(b+I),
\end{aligned}
\]
以及
\[
\begin{aligned}
\bar\varphi\bigl((a+I)(b+I)\bigr)
&=\bar\varphi(ab+I)\\
&=\varphi(ab)\\
&=\varphi(a)\varphi(b)\\
&=\bar\varphi(a+I)\bar\varphi(b+I).
\end{aligned}
\]
因此$\bar\varphi$为环同构。特别地，若$\varphi$为满射，则$\varphi(R)=R'$，故$R/I\cong R'$。
\end{proof}
\section{特征}
\begin{definition}
域$F$的单位元在加法群中若有有限阶$p$，称$F$的\textbf{特征}为$p$；若该阶无限，称特征为$0$。
记作$\operatorname{char}F$（或记为$\chi(F)$）。
也就是说，特征为正时是使$p1=0$成立的最小正整数。
\end{definition}

\begin{theorem}
域的特征只能为$0$或素数。
\end{theorem}
\begin{proof}
若1在加法群中有无限阶，按定义域的特征为0。

设1的加法阶为正整数$p$，即$p$是满足$p1=0$的最小正整数。
因为$1\neq0$，所以$p>1$。
假设$p$不是素数，则可写成
\[
p=uv,\qquad 1<u<p,\quad1<v<p.
\]
由$p$的最小性，
\[
u1\neq0,\qquad v1\neq0.
\]
但整数倍与乘法的相容性给出
\[
(u1)(v1)=(uv)1=p1=0.
\]
这与域无零因子矛盾。因此有限特征$p$必为素数，结论成立。
\end{proof}

\begin{proposition}
域$F$的所有非零元素在加法群中有相同的阶，等于$1$的加法阶。

特征为0时，对$n\in\Z\setminus\{0\},b\in F$，方程$nx=b$有唯一解
$x=(n1)^{-1}b$，常记为$\dfrac1n b$。

特征为$p>0$时，$n1=0$当且仅当$p\mid n$；此时方程$nx=b$仅在$b=0$时有解，并且每个$x$都是解。
若$p\nmid n$则仍有唯一解。
\end{proposition}
\begin{proof}
若$a\neq0$，则对任意整数$n$有
\[
na=n(a1)=a(n1).
\]
因此$n1=0$必推出$na=0$；反之，若$na=0$，因为$a$可逆，
\[
a^{-1}(na)=a^{-1}a(n1)=n1=0.
\]
故
\[
na=0\iff n1=0.
\]
所以$a$与1在加法群中满足相同的整数倍为零的条件，从而加法阶相同，包括两者均为无限阶的情形。

再讨论$nx=b$。由整数倍的性质，
\[
nx=(n1)x.
\]
若$\operatorname{char}F=0$且$n\neq0$，则$n1\neq0$，所以方程等价于
\[
x=(n1)^{-1}b.
\]
代回得到
\[
n\bigl((n1)^{-1}b\bigr)=(n1)(n1)^{-1}b=b,
\]
因此这确为唯一解。

\textit{补充正特征下的讨论。}若$\operatorname{char}F=p>0$，由于1的加法阶为$p$，
\[
n1=0\iff p\mid n.
\]
当$p\mid n$时，$nx=0$对所有$x$成立，所以$b=0$时每个$x$均为解，$b\neq0$时无解。
当$p\nmid n$时，$n1\neq0$，与前一情形相同，方程有唯一解$x=(n1)^{-1}b$。
\end{proof}

\begin{theorem}
\textbf{素域：}每个域含有唯一的最小子域。
特征为0时它同构于$\Q$；特征为素数$p$时它同构于$\F_p=\Z/p\Z$。
这些最小的域称\textbf{素域}。
\end{theorem}
\begin{proof}
先定义整数环到$F$的映射
\[
\eta:\Z\longrightarrow F,\qquad \eta(n)=n1.
\]
整数倍的运算法则给出
\[
\begin{aligned}
\eta(m+n)&=(m+n)1=m1+n1=\eta(m)+\eta(n),\\
\eta(mn)&=(mn)1=(m1)(n1)=\eta(m)\eta(n).
\end{aligned}
\]
所以$\eta$是环同态。

\textbf{特征为$p>0$的情形：}这时1的加法阶为$p$，故
\[
\ker\eta=\{n\in\Z:n1=0\}=p\Z.
\]
由环同态基本定理，
\[
\Z/p\Z\cong\eta(\Z)=\{0,1,2\cdot1,\ldots,(p-1)1\}.
\]
左边是域，因此右边的像也是域，并在$F$中构成一个子域，同构于$\F_p$。

\textbf{特征为0的情形：}此时$n1=0$只能在$n=0$时成立，所以$\eta$为单射，其像与整数环同构。
当$n\neq0$时，$n1\neq0$，在$F$中可取逆元，故把映射扩充到$\Q$：
\[
\widetilde\eta:\Q\longrightarrow F,\qquad
\widetilde\eta\left(\frac mn\right)=(n1)^{-1}(m1),\quad n\neq0.
\]

若$m/n=m'/n'$，则$mn'=m'n$，从而
\[
(m1)(n'1)=(m'1)(n1).
\]
在域$F$中乘以$(n1)^{-1}(n'1)^{-1}$，便得到
\[
(n1)^{-1}(m1)=(n'1)^{-1}(m'1).
\]
因此$\widetilde\eta$良定义。

记$M=m1,N=n1,M'=m'1,N'=n'1$，其中$N,N'$非零。
利用域乘法交换及分数通分，有
\[
\begin{aligned}
\widetilde\eta\left(\frac mn+\frac{m'}{n'}\right)
&=(NN')^{-1}(MN'+M'N)\\
&=N^{-1}M+(N')^{-1}M'\\
&=\widetilde\eta\left(\frac mn\right)
+\widetilde\eta\left(\frac{m'}{n'}\right),
\end{aligned}
\]
以及
\[
\begin{aligned}
\widetilde\eta\left(\frac mn\frac{m'}{n'}\right)
&=(NN')^{-1}MM'\\
&=(N^{-1}M)((N')^{-1}M')\\
&=\widetilde\eta\left(\frac mn\right)
\widetilde\eta\left(\frac{m'}{n'}\right).
\end{aligned}
\]
所以它是环同态。若$\widetilde\eta(m/n)=0$，则$m1=0$；因特征为0，有$m=0$，故核为$\{0\}$，映射单射。

设其像为$F_0$。它作为同态像已是子环，且包含$\widetilde\eta(1)=1$。
若$0\neq\widetilde\eta(m/n)\in F_0$，则$m\neq0$，并有
\[
\widetilde\eta\left(\frac mn\right)
\widetilde\eta\left(\frac nm\right)
=\widetilde\eta(1)=1.
\]
所以非零元素的逆仍在$F_0$，故$F_0$为子域，且$F_0\cong\Q$。

\textit{补充最小性与唯一性。}任意子域$K\subseteq F$都与$F$共用单位元1，因此包含1的所有整数倍。
正特征时，这说明$K$包含上面构造的$\eta(\Z)$。
特征为0时，$K$还包含每个非零$n1$的逆，因而包含全部$(n1)^{-1}(m1)$，即包含$F_0$。
所以构造出的子域包含于每个子域中，是最小子域。
若另有一个最小子域，两者必须互相包含，故相等，因此最小子域唯一。
\end{proof}



%\section{章末习题中的知识与解法}
%\noindent\textit{习题1--45。}
%% !TEX root = ../../main.tex
% 本文件可由 Part 2/chapter1.tex 输入。

\subsection*{群的公理、交换性与例子（习题1--16）}

\begin{problem}
\textbf{习题1：平方公式判别交换群。}若群$G$中任意$a,b$满足$(ab)^2=a^2b^2$，则$G$交换。
\end{problem}
\begin{proof}
从$abab=aabb$左消去$a$、右消去$b$，得$ba=ab$。
反向，交换群中该公式成立。
\end{proof}

\begin{problem}
\textbf{习题2：每个元素的平方为单位元。}若$a^2=e$对所有$a\in G$成立，则$G$交换。
\end{problem}
\begin{proof}
每个元素都等于自身的逆元，故$ab=(ab)^{-1}=b^{-1}a^{-1}=ba$。
\end{proof}

\begin{problem}
\textbf{习题3：有限消去半群是群。}非空有限集合上的结合运算若同时满足左右消去律，则构成群。
\end{problem}
\begin{proof}
固定$a$，左右乘$a$都是有限集的单射，所以均为满射。
可取$e,f$使$ae=a,fa=a$。
由$a(ex)=(ae)x=ax$左消去$a$得$ex=x$；由$(xf)a=x(fa)=xa$右消去$a$得$xf=x$。
于是$e=ef=f$，得到双边单位元。
满射性又给出$ab=e,ca=e$；$c=c(ab)=(ca)b=b$，故双边逆元存在。
有限性不能删去，例如正整数在加法下结合且可消去，却无单位元和逆元。
\end{proof}

\begin{problem}
\textbf{习题4：可解方程的结合系统。}非空集合上的结合运算，若对每对$a,b$，$ax=b$和$ya=b$都有解，则构成群。
\end{problem}
\begin{proof}
固定$a$，取$e,f$使$ae=a,fa=a$。
任意$x$可写成$ya$，所以$xe=y(ae)=ya=x$；任意$x$也可写成$az$，所以$fx=(fa)z=az=x$。
因此$f=fe=e$，存在双边单位元。
再取$ab=e,ca=e$，由$c=c(ab)=(ca)b=b$得双边逆元。
这里无需预设有限性或解的唯一性；构成群后，解的唯一性由消去律推出。
\end{proof}

\begin{problem}
\textbf{习题5：非交换平方公式的反例。}
\end{problem}
\begin{solution}
在$S_3$中取$x=(12),y=(23)$。则$x^2y^2=e$，而$xy=(123)$，$(xy)^2=(132)\neq e$。
\end{solution}

\begin{problem}
\textbf{习题6：任意偶数阶的非交换例子。}对$n>2$，构造一个$2n$阶非交换群。
\end{problem}
\begin{solution}
取正$n$边形的二面体群$D_n$。正文已经证明其阶为$2n$，且$ST=T^{-1}S\neq TS$。
\end{solution}

\begin{problem}
\textbf{习题7：反复共轭。}若$a^{-1}ba=b^r$，$r$为正整数，则对任意正整数$i$，$a^{-i}ba^i=b^{r^i}$。
\end{problem}
\begin{proof}
对$i$归纳。由$a^{-1}b^ma=(a^{-1}ba)^m=b^{rm}$，若结论对$i$成立，则
\[
a^{-(i+1)}ba^{i+1}=a^{-1}b^{r^i}a=b^{r^{i+1}}.
\]
\end{proof}

\begin{problem}
\textbf{习题8：取逆映射的同态性。}群$G$交换当且仅当$a\mapsto a^{-1}$是群同构。
\end{problem}
\begin{proof}
取逆映射总是双射。若它保乘法，则$a^{-1}b^{-1}=(ab)^{-1}=b^{-1}a^{-1}$，对所有$a,b$成立，故交换。
反向，交换群中的反序公式与保乘法公式相同。
\end{proof}

\begin{problem}
\textbf{习题9：由子集定义等价关系。}设$S\subseteq G$非空，定义$a\sim b\iff ab^{-1}\in S$。
该关系是等价关系当且仅当$S$为子群。
\end{problem}
\begin{proof}
若$S$为子群，则自反性由$e\in S$，对称性由取逆，传递性由
$(ab^{-1})(bc^{-1})=ac^{-1}$成立。
反之，自反性给出$e\in S$。若$u,v\in S$，则$u\sim e,e\sim v$（后者由$v\sim e$及对称性），传递性给出$uv^{-1}\in S$。
用子群判别法即可。
等价类是右陪集$Sb$，须留意$ab^{-1}$与$a^{-1}b$的次序。
\end{proof}

\begin{problem}
\textbf{习题10：$n\Z\cong\Z$。}这里$n$是正整数，讨论加法群同构。
\end{problem}
\begin{proof}
映射$k\mapsto nk$是$\Z\to n\Z$的保加法双射。
当$n>1$时它不是按通常乘法的环同构，因为一般$n(kl)\neq(nk)(nl)$。
\end{proof}

\begin{problem}
\textbf{习题11：克莱因四元群。}
$B=\{e,(12)(34),(13)(24),(14)(23)\}$为$S_4$的子群，且$B\cong U_4$不成立。
\end{problem}
\begin{proof}
记三个非单位元为$u,v,w$，直接计算得$u^2=v^2=w^2=e$，$uv=vu=w$，其余两两乘积同理。
故对积与逆封闭，构成子群，称克莱因四元群$V_4$。
$U_4=\{1,i,-1,-i\}$有4阶元素$i$，而$B$的非单位元均为2阶；同构保持阶，故不同构。
\end{proof}

\begin{problem}
\textbf{习题12：指数2的子群正规。}若$[G:H]=2$，则$H\trianglelefteq G$。
\end{problem}
\begin{proof}
当$a\in H$，$aH=H=Ha$。当$a\notin H$，左、右陪集都与$H$不交，且由陪集数为2，均等于$G\setminus H$。
故任意$a$均有$aH=Ha$。特别地，$|G|=2n$且$|H|=n$时结论成立。
\end{proof}

\begin{problem}
\textbf{习题13：偶数阶群有2阶元素。}
\end{problem}
\begin{proof}
把不等于自身逆元的元素配成$\{a,a^{-1}\}$。这些元素总数为偶数。
剩下的元素满足$a=a^{-1}$，其中已有$e$。总阶为偶数，剩余数也为偶数，因此至少还有一个$a\neq e$。
它满足$a^2=e$，阶为2。
\end{proof}

\begin{problem}
\textbf{习题14：二面体群的矩阵实现。}对$n\geq3$，令$\zeta=e^{2\pi i/n}$，
\[
A=\begin{pmatrix}0&1\\1&0\end{pmatrix},\qquad B=\begin{pmatrix}\zeta&0\\0&\zeta^{-1}\end{pmatrix}.
\]
则$\{B^k,AB^k:0\leq k<n\}\cong D_n$。
\end{problem}
\begin{proof}
$B^n=E,A^2=E,AB=B^{-1}A$。这些关系使乘积与逆仍有两种规定形状。
$B$的阶为$n$，故各$B^k$不同；各$AB^k$也不同。前一类为对角矩阵，后一类为非零反对角矩阵，两类不交。
所以该集合成$2n$阶群。
映射$T^k\mapsto B^k,ST^k\mapsto AB^k$是双射，并由相同的乘法关系保乘法，故为同构。
\end{proof}

\begin{problem}
\textbf{习题15：连续三个指数的乘积公式。}固定正整数$i$。若对所有$a,b\in G$及$k=i,i+1,i+2$均有$(ab)^k=a^kb^k$，则$G$交换。
\end{problem}
\begin{proof}
比较$(ab)^{i+1}=(ab)^i ab$与给定公式，得
$a^{i+1}b^{i+1}=a^ib^iab$。左消去$a^i$、右消去$b$，得$ab^i=b^ia$。
同理比较指数$i+1,i+2$，得$ab^{i+1}=b^{i+1}a$。
$a$于是与$(b^i)^{-1}b^{i+1}=b$交换。
\end{proof}

\begin{problem}
\textbf{习题16：有限阶元素的积可以无限阶。}在$\SL_2(\Q)$中令
\[
A=\begin{pmatrix}0&-1\\1&0\end{pmatrix},\qquad B=\begin{pmatrix}0&1\\-1&-1\end{pmatrix}.
\]
则$A$为4阶，$B$为3阶，$AB$为无限阶。
\end{problem}
\begin{proof}
直接计算$A^2=-E,A^4=E$，$B^2=\left(\begin{smallmatrix}-1&-1\\1&0\end{smallmatrix}\right),B^3=E$，且更小正幂不为$E$。
又$AB=\left(\begin{smallmatrix}1&1\\0&1\end{smallmatrix}\right)$，归纳得
$(AB)^m=\left(\begin{smallmatrix}1&m\\0&1\end{smallmatrix}\right)$，任意正整数$m$均不为$E$。
\end{proof}

\subsection*{子群、乘积与正规性（习题17--26）}

\begin{problem}
\textbf{习题17：交换群的有限阶元素成子群。}
\end{problem}
\begin{proof}
集合含$e$。若$a^m=e,b^n=e$，由交换性得$(ab^{-1})^{mn}=a^{mn}b^{-mn}=e$，故满足子群判别法。
交换性不能删去，习题16中的$A,B$均有限阶而$AB$无限阶。
\end{proof}

\begin{problem}
\textbf{习题18：只有有限多个子群的群必有限。}
\end{problem}
\begin{proof}
若有无限阶$a$，则$\langle a\rangle,\langle a^2\rangle,\langle a^3\rangle,\ldots$两两不同，因为无限循环群中的整数指数子群$m\Z$两两不同。
这与子群数有限矛盾。故所有元素有限阶。
每个元素属于一个有限循环子群，而循环子群也只有有限多个，因此$G$是有限多个有限集的并。
\end{proof}

\begin{problem}
\textbf{习题19：$D_n$的全部正规子群。}使用正文中$|D_n|=2n$的约定，$n\geq3$。
\end{problem}
\begin{solution}
所有$\langle T^d\rangle$（$d\mid n$）均正规，包含平凡子群$\langle T^n\rangle$。
此外有整个$D_n$；若$n$为偶数，还有
\[
\langle T^2,S\rangle,\qquad\langle T^2,ST\rangle.
\]
\end{solution}
\begin{proof}
循环群$\langle T\rangle$的子群由最小正指数与带余除法得到，均为$\langle T^d\rangle$，$d\mid n$。
反射把$T$共轭为$T^{-1}$，所以每个旋转子群均正规。
若子群$H$含反射$ST^k$，设$H\cap\langle T\rangle=\langle T^d\rangle$。
两反射之积为旋转，故$H=\langle T^d,ST^k\rangle$，其中$k$可模$d$选择。
若$H$正规，共轭$T(ST^k)T^{-1}=ST^{k-2}$必须仍在$H$，所以$d\mid2$。
因此$d=1$给出整个群，$d=2$要求$n$偶，给出$k=0,1$两群。
这两群的指数均为2，因而正规，列表完整。
\end{proof}

\begin{problem}
\textbf{习题20：子群乘积的判别。}对$H,K\leq G$，$HK$为子群当且仅当$HK=KH$。
\end{problem}
\begin{proof}
若$HK$为子群，则$HK=(HK)^{-1}=KH$。
若$HK=KH$，则$(HK)(HK)=H(KH)K=H(HK)K=HK$，且$(HK)^{-1}=KH=HK$。
它非空并对积、逆封闭，故为子群。
\end{proof}

\begin{problem}
\textbf{习题21：有限子群乘积的计数。}不要求$HK$为子群，也有
\[
|HK|=\frac{|H||K|}{|H\cap K|}.
\]
\end{problem}
\begin{proof}
映射$H\times K\to HK$，$(h,k)\mapsto hk$满。
若$h'k'=hk$，则$t=h^{-1}h'=kk'^{-1}\in H\cap K$，而$(h',k')=(ht,t^{-1}k)$。
反之每个$t\in H\cap K$都给出一个这样的表示，且彼此不同。
因此每个像有$|H\cap K|$个原像，得到公式。
\end{proof}

\begin{problem}
\textbf{习题22：正规子群的乘积。}若$M,N\trianglelefteq G$，则$MN=NM$且$MN\trianglelefteq G$。
若$M\cap N=\{e\}$，则$MN/N\cong M$。
\end{problem}
\begin{proof}
由$M$正规，$mn=n(n^{-1}mn)\in NM$，反向同理，故相等并成子群。
对$a\in G$有$a(MN)a^{-1}=(aMa^{-1})(aNa^{-1})=MN$，故正规。
定义$M\to MN/N$，$m\mapsto mN$，为满同态（$mnN=mN$），核为$M\cap N$。
第一同构定理给出更一般的$M/(M\cap N)\cong MN/N$，在交平凡时即得原结论。
\end{proof}

\begin{problem}
\textbf{习题23：中心化子与正规化子。}对非空$S\subseteq G$，定义
\[
C_G(S)=\{x:xa=ax\text{对所有 }a\in S\},\quad
N_G(S)=\{x:x^{-1}Sx=S\}.
\]
则两者均为子群，且$C_G(S)\trianglelefteq N_G(S)$。
\end{problem}
\begin{proof}
恒等元属于两者；与每个$a$交换的元素的积和逆也与$a$交换，所以$C_G(S)$为子群。
若$x,y$正规化$S$，则$(xy)^{-1}S(xy)=y^{-1}(x^{-1}Sx)y=S$；取逆同样保持$S$，所以$N_G(S)$为子群。
显然$C_G(S)\leq N_G(S)$。
对$n\in N_G(S),c\in C_G(S),a\in S$，$n^{-1}an\in S$，故
\[
(ncn^{-1})a=n c(n^{-1}an)n^{-1}
=n(n^{-1}an)c n^{-1}=a(ncn^{-1}).
\]
因此共轭仍在中心化子中，由正规性判别得结论。
\end{proof}

\begin{problem}
\textbf{习题24--26：2阶、4阶与素数阶群。}
\end{problem}
\begin{solution}
\textbf{24.}2阶群的非单位元$a$必为2阶，故$\{e,a\}\cong\{1,-1\}$。

\textbf{25.}4阶群恰有两种同构类型：循环群$\Z/4\Z$与$V_4$。
若有4阶元素，它生成全群；否则所有非单位元均为2阶，由习题2知交换。
取不同非单位元$a,b$，则$e,a,b,ab$两两不同，乘法表由$a^2=b^2=e,ab=ba$确定，正是$V_4$。
两者由是否有4阶元素区分。

\textbf{26.}若$|G|=p$为素数，任意$a\neq e$的阶整除$p$且不是1，故为$p$，生成$G$。
映射$\bar k\mapsto a^k$给出$\Z/p\Z\cong G$，所以任意两个$p$阶群同构。
\end{solution}

\subsection*{环的构造与单侧单位元（习题27--31）}

\begin{problem}
\textbf{习题27：整数对上的环。}在$\Z\times\Z$上定义
\[
(a,b)+(c,d)=(a+c,b+d),\qquad(a,b)(c,d)=(ac+bd,ad+bc).
\]
\end{problem}
\begin{solution}
映射$(a,b)\mapsto\left(\begin{smallmatrix}a&b\\b&a\end{smallmatrix}\right)$是到$M_2(\Z)$一个子环的保两种运算双射。
故环公理由矩阵环继承。该环交换，零元为$(0,0)$，单位元为$(1,0)$。
它有零因子，因为$(1,1)(1,-1)=(0,0)$，不是整环。
\end{solution}

\begin{problem}
\textbf{习题28：交换整数加法与乘法的角色。}令$a\oplus b=ab$，$a\odot b=a+b$，这是否使$\Z$成环？
\end{problem}
\begin{solution}
不成环。新“加法”是通常乘法，其单位元为整数1，但整数0没有新加法逆元，因为$0\oplus x=0$不可能等于1。
所以首先就不满足加法群公理。
\end{solution}

\begin{problem}
\textbf{习题29：改变坐标得到同构环。}在交换幺环$R$上令
$a\oplus b=a+b-1$，$a\odot b=a+b-ab$。
\end{problem}
\begin{proof}
双射$f(a)=1-a$满足
\[
f(a\oplus b)=f(a)+f(b),\qquad f(a\odot b)=f(a)f(b).
\]
所以新运算使$R$成一个与原环同构的交换幺环。
新加法零元为原元素1，新乘法单位元为原元素0，新加法逆元为$2-a$。
\end{proof}

\begin{problem}
\textbf{习题30：子环与母环的单位元和交换性。}
\end{problem}
\begin{solution}
\begin{enumerate}
\item 母环有$1$、子环无$1$：$R=\Z,S=2\Z$。
\item 母环无$1$、子环有$1$：$R=2\Z\times\Z,S=\{0\}\times\Z$，逐分量运算。
若$R$有$1=(u,v)$，首分量须满足$2u=2$，即$u=1\notin2\Z$；$S$的单位元却是$(0,1)$。
\item 两者均有单位元但不同：$R=M_2(\R)$，$S=\{\operatorname{diag}(a,0):a\in\R\}$，单位元分别是$E$与$\operatorname{diag}(1,0)$。
\item 母环不交换而子环交换：$R=M_2(\R)$，$S$为全部数量矩阵。
\end{enumerate}
\end{solution}

\begin{problem}
\textbf{习题31：单侧单位元。}若$e_La=a$对全部$a$成立，称$e_L$为左单位元；右单位元类似。
\end{problem}
\begin{proof}
\textbf{(i)}若同时有左单位元$e$、右单位元$f$，则$ef=f$且$ef=e$，故$e=f$为双边单位元。

\textbf{(ii)}若$R$无零因子且有左单位元$e$，对任意$a,b$，
$(ae-a)b=a(eb)-ab=0$。若环非零，固定非零$b$，无零因子得$ae=a$；零环的唯一元素本来就是双边单位元。

\textbf{(iii)}若有左单位元$e$却无右单位元，存在$a$使$d=a-ae\neq0$。
因为$d x=ax-a(ex)=0$对全部$x$成立，所以$(e+d)x=x$。
因此$e$与$e+d$是两个不同的左单位元。
\end{proof}

\subsection*{理想、幂零元与单位（习题32--42）}

\begin{problem}
\textbf{习题32：域没有非平凡双边理想。}
\end{problem}
\begin{proof}
理想$I$若非零，取$a\in I\setminus\{0\}$，则$a^{-1}a=1\in I$，故$I=F$。
同样证明适用于体。
\end{proof}

\begin{problem}
\textbf{习题33：$Ra$何时是理想。}若$R$交换，则$Ra=\{ra:r\in R\}$是双边理想；非交换时未必。
\end{problem}
\begin{proof}
差封闭：$ra-sa=(r-s)a$。左吸收：$t(ra)=(tr)a$；交换性给出右吸收$(ra)t=(rt)a$。
非交换反例：$R=M_2(\Q),a=E_{11}$。
$Ra$由第二列为0的矩阵组成，$E_{11}\in Ra$，却有$E_{11}E_{12}=E_{12}\notin Ra$。
注意：无单位元时$Ra$未必包含$a$，所以不能总把它当成“包含$a$的最小理想”。
\end{proof}

\begin{problem}
\textbf{习题34：理想的根。}在交换环$R$中定义
$\sqrt I=\{r\in R:r^n\in I\text{对某个正整数 }n\}$，则$\sqrt I$为理想。
\end{problem}
\begin{proof}
它含0。若$r^n,s^m\in I$，展开$(r-s)^{n+m-1}$，每个单项式不是含至少$n$个$r$，就是含至少$m$个$s$，因而在$I$中。
整数系数由加法子群封闭性处理，故$r-s\in\sqrt I$。
若$t\in R$，则$(tr)^n=t^nr^n\in I$，所以$tr\in\sqrt I$，交换性给出另一侧。
此处只用正整数幂；即使环无$1$，二项式展开中某一因子的次数为0时省略该因子，证明仍成立。
\end{proof}

\begin{problem}
\textbf{习题35：交换幺环仅有平凡理想时为域。}假设$1\neq0$。
\end{problem}
\begin{proof}
任意$a\neq0$，理想$Ra$含$a=1a$而非零，故为全环。
于是存在$r$使$ra=1$，交换性给出双边逆元，所以为域。
条件“至少两个元素”不能删，否则零环是反例。
\end{proof}

\begin{problem}
\textbf{习题36：全矩阵环是单环。}$M_n(\Q)$没有非平凡双边理想。
\end{problem}
\begin{proof}
若理想$I$非零，取$A=(a_{ij})\in I$及非零系数$a_{ij}$。
对任意$k,l$，$E_{ki}AE_{jl}=a_{ij}E_{kl}\in I$。
再左乘$a_{ij}^{-1}E_{kk}$得到$E_{kl}\in I$。
任意有理矩阵可写为$\sum q_{kl}E_{kl}$，而每个$q_{kl}E_{kl}=(q_{kl}E_{kk})E_{kl}\in I$，故$I$为全环。
只用理想吸收性，无需假设理想预先为$\Q$-线性空间。
\end{proof}

\begin{problem}
\textbf{习题37：$aba=0$判别零因子。}若$a,b\neq0$且$aba=0$，则$a$为左零因子或右零因子。
\end{problem}
\begin{proof}
若$ba=0$，非零$b$使$a$为右零因子；若$ba\neq0$，则$a(ba)=0$，使$a$为左零因子。
\end{proof}

\begin{problem}
\textbf{习题38：$1-a$的有限几何级数逆。}幺环中若$a^n=0$，则$a$称\textbf{幂零元}，且$1-a$可逆。
\end{problem}
\begin{proof}
令$u=1+a+\cdots+a^{n-1}$。望远镜相消得到$(1-a)u=u(1-a)=1-a^n=1$。
无须整个环交换，因为只涉及同一元素的幂。
\end{proof}

\begin{problem}
\textbf{习题39：交换环的幂零元构成理想。}
\end{problem}
\begin{proof}
这是习题34中$I=\{0\}$的情形。
若$a^n=b^m=0$，则$(a-b)^{n+m-1}=0$；若$r\in R$，则$(ra)^n=r^na^n=0$。
非交换环不成立，例如$M_2(\Q)$中$E_{12},E_{21}$均幂零，但$(E_{12}+E_{21})^2=E$。
\end{proof}

\begin{problem}
\textbf{习题40：有限幺环的单侧逆也是双侧逆。}若$R$有限且$xy=1$，则$yx=1$。
\end{problem}
\begin{proof}
记$L_x(r)=xr,L_y(r)=yr$，则$L_xL_y=\operatorname{id}$，所以$L_x$满。
有限集上的满射为双射，因此$L_y=L_x^{-1}$，有$L_yL_x=\operatorname{id}$。
代入$r=1$得到$yx=1$。
\end{proof}

\begin{problem}
\textbf{习题41：单侧逆不成双侧逆时有无穷多个右逆。}若幺环中$ab=1$但$ba\neq1$，则满足$ax=1$的$x$有无穷多个。
\end{problem}
\begin{proof}
令$c=1-ba\neq0$，则$ac=0,cb=0$。
对每个$n\geq0$，令$x_n=b+ca^n$，有$ax_n=1$。
若$x_i=x_j$，$i<j$，则$ca^i=ca^j$。
右乘$b^j$，利用$ab=1$给出的$a^jb^j=1$及$a^ib^j=b^{j-i}$，得
$0=cb^{j-i}=c$，矛盾。故这些右逆两两不同。
\end{proof}

\begin{problem}
\textbf{习题42：唯一的广义逆推出体。}环$R$至少有两个元素，若每个$a\neq0$恰有一个$b$满足$aba=a$，则：无零因子，$bab=b$，有单位元，且$R$为体。
\end{problem}
\begin{proof}
若$a\neq0$且$ax=0$，则$a(b+x)a=a$，由$b$的唯一性得$x=0$；若$xa=0$也同理。
故无左右零因子。

固定$a\neq0$及其$b$，由$aba=a$对任意$r\in R$有
$a(bar-r)=0$，所以$bar=r$；又$(rab-r)a=0$，所以$rab=r$。
故$ba$为左单位元、$ab$为右单位元，习题31(i)给出$ba=ab=1$。
于是$b\neq0$，$bab=b$；每个非零$a$有双边逆元$b$，所以$R$为体。
\end{proof}

\subsection*{函数环、有限域与体中的恒等式（习题43--45）}

\begin{problem}
\textbf{习题43：连续函数环。}对$C[0,1]$，每个真理想$I\neq C[0,1]$有公共零点；对非零函数$f\in C[0,1]$，$f$为零因子当且仅当其零点集包含一个非空的相对开区间。
\end{problem}
\begin{proof}
\textbf{(i)}若无公共零点，每个$t\in[0,1]$都有$f_t\in I$使$f_t(t)\neq0$。
连续性给出这些函数非零的开邻域，它们覆盖$[0,1]$；紧致性给出有限子覆盖，设对应函数为$f_1,\ldots,f_m$。
$g=\sum f_j^2\in I$且处处正，故$1/g$连续，理想吸收性给出$1=g(1/g)\in I$，与$I$真矛盾。

\textbf{(ii)}对非零$f$，若$fg=0$且$g\neq0$，选$g(t)\neq0$的一点。
连续性使$g$在一个非空相对开区间上非零，故$f$在该区间为0。
反之，若$f$在一个非空相对开区间为0，可在其中选一个非零连续三角形状函数$g$，并使$g$在区间外为0，则$fg=0$。
零函数也满足零点集条件，但按零因子的定义，它不是零因子。因此这里须限定$f\neq0$。
\end{proof}

\begin{problem}
\textbf{习题44：有限域上的矩阵计数。}若$F=\F_p$，则
\[
|M_n(F)|=p^{n^2},\qquad
|\GL_n(F)|=\prod_{k=0}^{n-1}(p^n-p^k).
\]
\end{problem}
\begin{proof}
矩阵有$n^2$个独立位置，各有$p$种选择，故第一式成立。
可逆矩阵的列向量必须构成有序基。
第一列有$p^n-1$种；已选$k$个线性无关列后，它们张成的空间有$p^k$个向量，下一列必须在其外，有$p^n-p^k$种。
乘法原理即得第二式。$n=1$时为$p-1$，与$F^\times$一致。
\end{proof}

\begin{problem}
\textbf{习题45：华罗庚恒等式。}在体$K$中，若$a,b\neq0$且$ab\neq1$，则
\[
a-\bigl(a^{-1}+(b^{-1}-a)^{-1}\bigr)^{-1}=aba.
\]
\end{problem}
\begin{proof}
先验证所有逆元都有定义。令$c=b^{-1}-a$，若$c=0$则$ab=1$，矛盾，所以$c\neq0$。
令$d=cba=a-aba$，因体无零因子且$c,b,a$均非零，$d\neq0$。
计算时保留乘法次序：
\[
\bigl(a^{-1}+c^{-1}\bigr)d
=a^{-1}(a-aba)+c^{-1}(cba)
=(1-ba)+ba=1.
\]
因此$a^{-1}+c^{-1}=d^{-1}$，其逆为$d$，于是左边为$a-d=aba$。
对$x\neq0,1$，还可验证$(x^{-1}-1)^{-1}=(1-x)^{-1}-1$：
$x^{-1}-1=x^{-1}(1-x)$，所以其逆为$(1-x)^{-1}x=(1-x)^{-1}-1$。
上述直接证明同时处理了非交换乘法与逆元存在条件。
\end{proof}

%
%\section*{复习时应核对的条件}
%\addcontentsline{toc}{section}{复习时应核对的条件}
%\begin{enumerate}
%\item 验证群：先说明运算封闭，再说明结合律、单位元、逆元；单侧公理须先推出双侧性质。
%\item 验证子群或子环：不能忘记非空；分别用$ab^{-1}$或$a-b,ab$判别。
%\item 非交换乘法不能随意改变次序；$(ab)^{-1}=b^{-1}a^{-1}$。
%\item 陪集不是子群；$G/H$成商群前先核对$H$正规。
%\item 同态的反像是集合；核描述“被压成单位元或零元”的部分，像描述得到的结构。
%\item 环不必有$1$；子环可以没有$1$，也可能有与母环不同的$1$；同构与一般同态在保持$1$上有区别。
%\item 无零因子才有非零消去律；一侧逆元与双侧逆元要区分。
%\item 理想必须被母环元素两侧吸收；商环必须证明代表选择无关。
%\item 有限整环是域，不能遗漏有限性、交换性、单位元和$1\neq0$。
%\item 特征为0不表示$1$的加法阶为0，而表示该阶无限；特征为$p$时只有$p\nmid n$才允许除以$n$。
%\end{enumerate}
